% ===================================================================
%  TABELO N-RO 14 — La strato. — La komercistoj.
%  Traduction 2026 d'apres texto/io, controlee sur texto/fr et
%  texto/es. Consigne : voir texto/eo/10-tabelo-01.tex.
% ===================================================================

%%K t14-apar-1 apar
\VUsaut{4.0mm}
\VUtitre{15.0pt}{60}{TABELO N\textsuperscript{o} 14}
\VUsaut{3.0mm}

%%K t14-tit-1 sub
\VUpk{11.4pt}{40}{La strato. --- La komercistoj.}
\VUsaut{3.0mm}

%%K t14-01-1 p
1. --- Mia amiko Johano, kiu loĝas la kamparon, venis pasigi tri tagojn en mia
familio. \VUgras{Ĝis} tiam li ne estis havinta la okazon vidi grandan urbon; pro
tio ni uzas ĉiujn niajn tagojn promenante, kaj mi klarigas al li tion kion li
ne konas.

%%K t14-02-1 p
2. --- Unue ni iris viziti la \VUgras{observatorion} \textsuperscript{(1)}, kiu
tre interesis lin. Revenante tra la \VUgras{strato} \textsuperscript{(3)} kie
estas la \VUgras{turka banejo} \textsuperscript{(2)}, ni renkontis
\VUgras{funebron} \textsuperscript{(4)}. Antaŭ la \VUgras{funebra konvojo}
marŝis la \VUgras{klerikaro}, antaŭiranta la \VUgras{mortintveturilon}, en kiu
estis metita la \VUgras{ĉerko}. Multaj amikoj akompanis la \VUgras{funebrantan}
familion.

%%K t14-02-2 p
Post la paso de la sekvantaro, ni transiris la straton antaŭ la granda
\VUgras{bierejo} \textsuperscript{(5)} en kiu multaj \VUgras{konsumantoj} trinkis
\VUgras{bieron} sur la teraso, ŝirmate de la vitrita \VUgras{markezo}
\textsuperscript{(6)}.

%%K t14-03-1 p
3. --- Johano estis mirigita de la \VUgras{ŝildo} metita inter la magazeno de la
\VUgras{likvoristo} \textsuperscript{(7)} kaj tiu de la \VUgras{muzikvendisto}
\textsuperscript{(10)}. Li demandis min pri ĝia signifo; mi respondis ke ĝi
montras la \VUgras{konsulejon} \textsuperscript{(8)} anglan, kaj rimarkigis al
li la \VUgras{konsulon} \textsuperscript{(9)} kiu estas sur la balkono.

%%K t14-tit-2 sub
\VUpk{10.2pt}{40}{Esploro al butikoj.}
\VUsaut{3.0mm}

%%K t14-04-1 p
4. --- Kompreneble ni haltis antaŭ la elmetejo de la \VUgras{muzikiloj},
\VUgras{harmoniumoj}, \VUgras{orgenoj} kaj \VUgras{pianoj}; ni rigardis la
ilustritajn kovrilojn de la \VUgras{partituroj} kaj de la \VUgras{muzikaĵoj} por
piano, kaj admiris la \VUgras{liron} el orita ligno, kiu uzatas kiel ŝildo.

%%K t14-05-1 p
5. --- La polvaj nuboj faritaj de la \VUgras{balaistoj} \textsuperscript{(11)}
kaj la \VUgras{balaveturilo} \textsuperscript{(12)} devigis nin rapide pasi
antaŭ la belaj \VUgras{meblaroj} de la \VUgras{meblovendisto}
\textsuperscript{(13)} kaj antaŭ la elmetejo de \VUgras{fornoj} kaj
\VUgras{lampoj} de la \VUgras{lad(if)isto} \textsuperscript{(14)}.

%%K t14-06-1 p
6. --- Cetere, ĉi-loke, ni estis devigitaj forlasi la trotuaron, ĉar ĝi estis
barita per pakoj, kiujn oni prenis el \VUgras{transportveturilo}
\textsuperscript{(16)} por loki ilin en \VUgras{magazenon}
\textsuperscript{(15)} kiun oni ĵus estis luinta.

%%K t14-06-2 p
Tio malhelpis nin vidi la belajn veturilojn de la \VUgras{veturilfaristo}
\textsuperscript{(17)}, sed ne malhelpis nin halti por rigardi la
\VUgras{pakiganton} \textsuperscript{(19)} farantan keston por sia najbaro, la
\VUgras{vendisto de bicikloj kaj aŭtomobiloj} \textsuperscript{(18)}. Poste,
ĉar jam estis iom malfrue, ni revenis al la hejmo.

%%K t14-tit-3 sub
\VUpk{10.2pt}{40}{Promeno tra la urbo.}
\VUsaut{3.0mm}

%%K t14-07-1 p
7. --- En la sekva tago, mia patrino proponis al Johano ke oni fotografu lin,
poste ŝi komisiis min por akompani lin al la \VUgras{fotografisto}
\textsuperscript{(20)}. Post la tagmanĝo, ni iris en la \VUgras{laborejon} de la
artisto, ĉe kiu ni estis devigitaj sufiĉe longe atendi. Por okupi nin, ni
rigardis la \VUgras{portretojn} \textsuperscript{(22)} pendantajn ĉe la muro.

%%K t14-07-2 p
Kiam estis nia vico, la laboro ne daŭris longe, kaj tuj kiam mia amiko finis
pozi antaŭ la \VUgras{fotografilo} \textsuperscript{(21)}, ni malsupreniris.

%%K t14-08-1 p
8. --- Sur la unua etaĝo, ni vidis tra la pordo de la \VUgras{frizisto}
\textsuperscript{(23)}, kiu estas malfermita, lian metilernanton kiu tondis la
hararon de kliento.

%%K t14-08-2 p
Alveninte en la straton, ni pasis antaŭ la \VUgras{tabakvendejo}
\textsuperscript{(24)}. Johano estis tiel amuzata de la ruĝa \VUgras{lanterno}
\textsuperscript{(25)} kaj la \VUgras{dika pipo} kiu estas uzata kiel
\VUgras{ŝildo} \textsuperscript{(26)}, ke li kolizias kun du maljunaj sinjoroj
kiuj konversaciis \VUgras{flarante tabakon} \textsuperscript{(27)}.

%%K t14-tit-4 sub
\VUpk{10.2pt}{40}{La kukejo.}
\VUsaut{3.0mm}

%%K t14-09-1 p
9. --- La ĉiulandaj \VUgras{bankbiletoj} kaj \VUgras{moneroj} elmetitaj en la
\VUgras{vitrino} de la \VUgras{monŝanĝisto} \textsuperscript{(29)} kelkamomente
altiris nian atenton, sed ĉar estis la horo manĝeti, ni eniris ĉe la
\VUgras{kukiston} \textsuperscript{(31)} sen halti pli longe.

%%K t14-10-1 p
10. --- Vidante ĉiujn \VUgras{frandaĵojn} kiujn enhavis la magazeno,
\VUgras{kukojn, mielkukojn, biskvitojn} kaj ĉiaspecajn \VUgras{bombonojn}, ni
ne facile povis elekti. Fine Johano elektis la \VUgras{kremkukojn}, kaj mi
prenis nur malgrandan \VUgras{ĉeriztarton}, ĉar la sukeraĵoj \VUgras{dolorigas
miajn dentojn}, kaj mi tute ne deziras iri tro ofte viziti la \VUgras{dentiston}
\textsuperscript{(30)}.

%%K t14-tit-5 sub
\VUpk{10.2pt}{40}{Maŝinskribado.}
\VUsaut{3.0mm}

%%K t14-11-1 p
11. --- Por montri al mia amiko \VUgras{skribmaŝinon} pri kiu li estis aŭdinta
paroli, sed kiun li ne konis, ni eniris la \VUgras{skribĉambron}
\textsuperscript{(32)} de mia \VUgras{kuzo} \textsuperscript{(34)}. Ni trovis
lin diktantan siajn leterojn al sia \VUgras{sekretario} \textsuperscript{(35)}
kiu estas \VUgras{stenografisto}. Apenaŭ kiam letero estis finita,
\VUgras{maŝinskribisto} \textsuperscript{(33)} rekopiis ĝin per sia maŝino.
Dezirante ne interrompi la okupojn de mia parenco, ni restis ĉe li nur dum
kelkaj momentoj.

%%K t14-12-1 p
12. --- Sub la oficejo de mia kuzo, loĝas granda \VUgras{horloĝist-juvelisto}
\textsuperscript{(37)} kiu estas krome orlaboristo. Lia belega magazeno enhavas
multajn \VUgras{horloĝojn, pendolhorloĝojn, poŝhorloĝojn}, \VUgras{ĉenetojn}
kaj multekostajn \VUgras{juvelojn}, ekzemple \VUgras{perlokolĉenojn}, orajn
\VUgras{braceletojn, orelringojn} kaj \VUgras{fingroringojn} ornamitajn per
\VUgras{gemoj} kaj \VUgras{diamantoj}. Unu el la vitrinoj estas aparte
rezervita por la \VUgras{oraĵoj} kaj enhavas gravan kolektaĵon da arĝentaj
\VUgras{manĝilaroj}.

%%K t14-13-1 p
13. --- Turnante ĉe la angulo de la strato, mi rimarkis ke oni estis ŝanĝinta
la \VUgras{plakon} \textsuperscript{(36)} metitan apud la \VUgras{numero} de la
domo. Mi estis laŭte dironta mian rimarkon, kiam la ĝojaj sonoj de \VUgras{la}
militista muziko estis aŭdataj. Ni ekkuris renkonte al la regimento, sen
pripensi ke \VUgras{ĝi} estis pasonta antaŭ la magazenoj de la
\VUgras{ĉapelisto} \textsuperscript{(39)} kaj de la \VUgras{gantisto}
\textsuperscript{(40)}, kaj ke ni estus tiel same bone lokitaj por vidi ĝian
defilon restante antaŭ la vitrino de la \VUgras{optikisto}
\textsuperscript{(38)}, tute provizita per \VUgras{nazbinokloj, orelbinokloj}
kaj \VUgras{lornetoj}.

%%K t14-tit-6 sub
\VUpk{10.2pt}{40}{Marŝdefilo.}
\VUsaut{3.0mm}

%%K t14-14-1 p
14. --- Antaŭ la \VUgras{regimento} \textsuperscript{(41)} marŝis la
\VUgras{tamburestro} \textsuperscript{(42)} sekvata de la \VUgras{tamburistoj}
\textsuperscript{(43)}, la \VUgras{klarionistoj} \textsuperscript{(44)} kaj la
tuta \VUgras{muzikistaro}. La \VUgras{generalo} \textsuperscript{(45)} kiu estis
sur la placo kun sia adjutanto, estis haltinta apud la \VUgras{akvoŝprucilo}
\textsuperscript{(49)} de la \VUgras{fontano} \textsuperscript{(48)}, por
ĉeesti la \VUgras{defilon}.

%%K t14-14-2 p
\VUgras{La stabaj oficiroj} \textsuperscript{(46)}, la \VUgras{kolonelo} kaj la
\VUgras{leŭtenantkolonelo} salutis lin per spado. \VUgras{La majoroj} estis
antaŭ siaj \VUgras{batalionoj} \textsuperscript{(47)} kaj ĉiu \VUgras{kapitano}
komandis sian \VUgras{kompanion}, dum la \VUgras{leŭtenantoj},
\VUgras{subleŭtenantoj} kaj \VUgras{suboficiroj} okupis sian kutiman lokon
apud siaj viroj.

%%K t14-15-1 p
15. --- Multaj pasantoj estis grupiĝintaj sur \VUgras{la vojkruco}
\textsuperscript{(50)}; la komercistoj, la \VUgras{ombrelisto}
\textsuperscript{(51)}, la \VUgras{tinkturisto} \textsuperscript{(53)}, la
\VUgras{armilfaristo} \textsuperscript{(54)} eliris por vidi la
\VUgras{soldatojn}. Ĉiuj veturiloj, \VUgras{fiakroj} \textsuperscript{(52)},
\VUgras{mastroveturiloj} \textsuperscript{(57)} kaj ankaŭ la
\VUgras{akvumfudro} \textsuperscript{(56)}, gariĝis laŭlonge de la trotuaro.

%%K t14-tit-7 sub
\VUpk{10.2pt}{40}{Scenoj ĉe la strato.}
\VUsaut{3.0mm}

%%K t14-15-2 p
\VUgras{Komercvojaĝisto} \textsuperscript{(55)} portanta per la mano siajn
\VUgras{specimenskatolojn}, rapide transiris la straton, kaj la personoj
sidantaj sur la \VUgras{imperialo} \textsuperscript{(59)} de la
\VUgras{omnibuso} \textsuperscript{(58)} returnis sin por ĝui la vidindaĵon.
Kompatinda \VUgras{kantisto vaganta} \textsuperscript{(60)} kun sia
\VUgras{turnorgeno} \textsuperscript{(61)} devis interrompi sian
\VUgras{kanzonon}, ĉar la bruado de la tamburoj kovris lian voĉon.

%%K t14-16-1 p
16. --- Kiam la regimento alvenis antaŭ la \VUgras{ŝtofkomercejon}
\textsuperscript{(64)} la \VUgras{kudristinoj} \textsuperscript{(63)} kaj la
\VUgras{tajloroj} \textsuperscript{(62)} forlasis siajn \VUgras{kudrmaŝinojn}
kaj sian laboron por rigardi tra la fenestro. \VUgras{La komercisto}
\textsuperscript{(65)}, kiu ĵus estis metinta novajn \VUgras{kostumojn}
\textsuperscript{(66)} ĉe sia \VUgras{elmetejo}, devis metigi interne la
\VUgras{drapaĵojn} \textsuperscript{(67)} kaj la \VUgras{tolaĵojn} instalitajn
sur la trotuaro, apud la \VUgras{kloakaperturo} \textsuperscript{(68)}, time ke
la homamaso renversos ilin, eĉ maljuna \VUgras{kolportisto}
\textsuperscript{(69)}, de kiu pordistino marĉandis ŝtrumpojn, estis devigita
eniri koridoron por fini sian vendon.

%%K t14-16-2 p
Ni akompanis la regimenton ĝis la kazerno, \VUgras{kaj}, tre kontentaj pri nia
tago, ni revenis je la horo de la vespermanĝo.

%%K t14-tit-8 sub
\VUpk{10.2pt}{40}{Diversaj komercoj kaj profesioj.}
\VUsaut{3.0mm}

%%K t14-17-1 p
17. --- En la sekvanta tago, mia patro proponis al ni viziti la presejon de la
ĉi-tiea ĵurnalo. Ni entuziasme akceptis.

%%K t14-17-2 p
La \VUgras{presisto} estas ankaŭ \VUgras{libristo} \textsuperscript{(70)}
\VUgras{(= librovendisto)}, \VUgras{eldonisto}, \VUgras{papervendisto} kaj
\VUgras{bindisto}. Li instalis sur la unua etaĝo la \VUgras{redaktejon}
\textsuperscript{(71)} de la ĵurnalo, kaj ankaŭ la \VUgras{gravurejon}, en kiu
laboras la \VUgras{litografistoj} \textsuperscript{(72)} kaj la
\VUgras{kuprogravuristoj} \textsuperscript{(73)}, fine la
\VUgras{kompostlaborejon}, en kiu estas la \VUgras{preslaboristoj}
\textsuperscript{(74)}. La \VUgras{presejo} \textsuperscript{(75)}
proprasence dirita, la \VUgras{prespresiloj} kaj la diversaj maŝinoj estas sur
la teretaĝo.

%%K t14-tit-9 sub
\VUpk{10.2pt}{40}{Johano faras tre multajn demandojn.}
\VUsaut{3.0mm}

%%K t14-18-1 p
18. --- Elirante el la presejo, Johano tre mirigita haltis antaŭ malgranda
vitrita skatolo metita sur kolonon, kontraŭ la \VUgras{galanteriejo}
\textsuperscript{(77)} kaj demandis por kio ĝi uzatas. Mia patro klarigis al li,
ke ĝi estas \VUgras{incendiavertilo} \textsuperscript{(76)}. Cetere ĉio, en la
urbo, estis por mia amiko, mirindaĵo, de la simpla \VUgras{ŝtonfontano}
\textsuperscript{(78)} kaj la malpeza \VUgras{portotriciklo}
\textsuperscript{(80)} kondukata de \VUgras{grumo} \textsuperscript{(79)} kun
livreo, ĝis la \VUgras{poŝtveturilo} \textsuperscript{(81)}.

%%K t14-19-1 p
19. --- Dum mia patro kelkamomente forlasis nin por paroli kun la
\VUgras{impostokolektisto} \textsuperscript{(83)}, kiu estis haltinta por
konversacii kun \VUgras{advokato} \textsuperscript{(82)} kaj \VUgras{notario}
\textsuperscript{(84)}, ni amuzis nin sekvante per la okuloj la cirkuladon sur
la strato. La plej diversaj personoj pasis antaŭ ni : \VUgras{kukista
metilernanto} \textsuperscript{(85)}, portanta en korbo \VUgras{pasteĉon}
belegan, venis malantaŭ \VUgras{vestovendisto} \textsuperscript{(86)};
\VUgras{lanternekbruligisto} \textsuperscript{(87)} estis konversacianta kun
\VUgras{gasisto} \textsuperscript{(88)}, \VUgras{monaĥino}
\textsuperscript{(89)} tenanta mane orfinon estis reiranta en sian monaĥejon.

%%K t14-tit-10 sub
\VUpk{10.2pt}{40}{Dum nia reveno hejmen.}
\VUsaut{3.0mm}

%%K t14-20-1 p
20. --- En la momento kiam mia patro rekuniĝis kun ni, Johano ridegis, vidante
la malpuran vizaĝon kaj la strangan vestaĉon de du \VUgras{kamenskrapistoj}
\textsuperscript{(91)} tute fulgokovritaj.

%%K t14-21-1 p
21. --- Mi volis aĉeti de la \VUgras{bukedistino} \textsuperscript{(92)} planton
por mia patrino, sed mia patro rimarkigis al mi, ke ĝi estus tre peza por
kunporti kaj ke estus preferinde akiri \VUgras{oranĝojn}. Ni alproksimiĝis al
la \VUgras{vendistino} \textsuperscript{(94)} kaj kiam la \VUgras{edukistino}
\textsuperscript{(93)}, kiu estis elektanta por sia edukato, finis sian aĉeton,
ni igis servi nin.

%%K t14-22-1 p
22. --- Revenante al la hejmo, ni renkontis plurajn konatojn : malnova amikino
de mia familio, kiu apogis la brakon sur sia \VUgras{kunula sinjorino}
\textsuperscript{(95)}, la \VUgras{inĝeniero} \textsuperscript{(96)} pri la
pontoj kaj vojoj, kaj unu el miaj kuzinoj kun sia \VUgras{infana gardistino}
\textsuperscript{(98)}.

%%K t14-22-2 p
Poste ni renkontis la \VUgras{ĉapelistinon} \textsuperscript{(97)} de mia
patrino kaj du \VUgras{monaĥojn} \textsuperscript{(99)} fremdajn al la urbo,
kiuj aliris nin por demandi mian patron pri la vojo al la stacidomo.
\VUsaut{6.0mm}
\VUfilet{22mm}
